\documentclass{article}

\usepackage[preprint]{aistats2027}
\usepackage[round,authoryear]{natbib}
\usepackage{amsmath,amssymb,amsthm}
\usepackage{mathtools}
\usepackage{microtype}
\usepackage{booktabs}
\usepackage{url}

\renewcommand{\bibname}{References}
\renewcommand{\bibsection}{\subsubsection*{\bibname}}
\bibliographystyle{apalike}

\newtheorem{theorem}{Theorem}
\newtheorem{proposition}[theorem]{Proposition}
\newtheorem{corollary}[theorem]{Corollary}
\newtheorem{lemma}[theorem]{Lemma}
\newcommand{\E}{\mathbb E}
\newcommand{\Prb}{\mathbb P}
\newcommand{\one}{\mathbf 1}
\newcommand{\TV}{d_{\rm TV}}
\newcommand{\supp}{\operatorname{supp}}
\newcommand{\ip}[2]{\langle #1,#2\rangle}

\begin{document}
\runningauthor{}

\onecolumn
\begin{center}
 {\Large\bfseries Supplement to ``The Confidence Cost of Pairwise
 Distribution Learning''}\\[0.75em]
 Pahan Dewasurendra
\end{center}

This supplement gives complete proofs of the winner-chain occupation bound,
the arbitrary-start upper bound, and the adaptive confidence lower bound. It
also proves the bounded-menu lower bound stated as an open-direction result
in the main paper and records the numerical and source audits. Numbering is
local to the supplement.

\section{Pairwise Winner Chain}
\label{app:chain}

Fix a distribution $p\in\Delta_n$. From state $i$, the chain chooses
$j\ne i$ uniformly, queries the pair $\{i,j\}$, and makes the returned label
the next state. If $p_i+p_j>0$, then
\begin{equation}
 P(i,j)=\frac{p_j}{(n-1)(p_i+p_j)},\qquad i\ne j.
 \label{eq:app-kernel}
\end{equation}
If $p_i=p_j=0$, the pair response may follow any fixed convention. The
diagonal is the remaining row mass.

\begin{lemma}[Stationarity, reversibility, and minorization]
\label{lem:app-basic}
The distribution $p$ is stationary. On its positive support the chain is
reversible, and for every state $i$,
\begin{equation}
 P(i,\cdot)\ge\frac1{n-1}p(\cdot)
 \label{eq:app-minorization}
\end{equation}
coordinatewise.
\end{lemma}

\begin{proof}
For distinct positive-mass states $i,j$,
\[
 p_iP(i,j)=\frac{p_ip_j}{(n-1)(p_i+p_j)}=p_jP(j,i).
\]
Transitions from a positive-mass state to a zero-mass state have probability
zero. Detailed balance on the support proves reversibility and stationarity.

For the minorization, first take $j\ne i$. If $p_i+p_j>0$, then
\[
 P(i,j)=\frac1{n-1}\frac{p_j}{p_i+p_j}
 \ge\frac{p_j}{n-1},
\]
since $p_i+p_j\le1$. If both masses vanish, the right side is zero. If
$p_i=0$, the diagonal inequality is immediate. If $p_i>0$, then against
every challenger $j$, the probability of remaining at $i$ is at least
$p_i/(p_i+p_j)\ge p_i$. Averaging over challengers gives
$P(i,i)\ge p_i\ge p_i/(n-1)$.
\end{proof}

For $n>2$, the minorization gives a Markov kernel $R$ such that
\begin{equation}
 P=\frac1{n-1}\Pi_p+\left(1-\frac1{n-1}\right)R,
 \label{eq:app-doeblin}
\end{equation}
where every row of $\Pi_p$ is $p$. Thus, for any initial law $\mu$,
\begin{equation}
 \TV(\mu P^B,p)
 \le\left(1-\frac1{n-1}\right)^B
 \le e^{-B/(n-1)}.
 \label{eq:app-mixing}
\end{equation}
For $n=2$, one query returns an exact $p$ draw from either state.

\section{Positive Spectrum and Coordinate Variance}
\label{app:variance}

We work on the positive support of $p$. Zero-mass coordinates contribute
nothing to TV under stationarity. Let $m=|\supp(p)|$. For a function
$f\in L_2(p)$, reversibility gives the Dirichlet form
\begin{equation}
 \ip{f}{(I-P)f}_p
 =\sum_{1\le i<j\le m}
 \frac{p_ip_j}{(n-1)(p_i+p_j)}(f_i-f_j)^2.
 \label{eq:app-dirichlet}
\end{equation}

\begin{lemma}[Positive semidefinite kernel]
\label{lem:app-psd}
Every eigenvalue of $P$ on $L_2(p)$ lies in $[0,1]$.
\end{lemma}

\begin{proof}
The Markov and reversibility properties put the spectrum in $[-1,1]$. For
each positive pair,
\begin{equation}
 \frac{p_ip_j}{p_i+p_j}(f_i-f_j)^2
 \le p_if_i^2+p_jf_j^2.
 \label{eq:app-edge-upper}
\end{equation}
Indeed, the right side minus the left side is
$(p_if_i+p_jf_j)^2/(p_i+p_j)$. Each coordinate appears in at most $m-1$
positive pairs. Summing \eqref{eq:app-edge-upper} and dividing by $n-1$
gives
\[
 \ip{f}{(I-P)f}_p\le\ip{f}{f}_p.
\]
Hence $\ip{f}{Pf}_p\ge0$ for every $f$.
\end{proof}

For coordinate $i$, put $f_i(x)=\one\{x=i\}-p_i$ and
\begin{equation}
 H_i=\ip{f_i}{(I-P)^{-1}f_i}_p,
 \label{eq:app-green}
\end{equation}
where the inverse acts on the mean-zero subspace.

\begin{lemma}[Finite-horizon coordinate variance]
\label{lem:app-coordinate-variance}
For a stationary length-$T$ trajectory,
\begin{equation}
 \operatorname{Var}_p\!\left(\frac1T\sum_{t=1}^Tf_i(X_t)\right)
 \le\frac{2H_i}{T}.
 \label{eq:app-coordinate-variance}
\end{equation}
\end{lemma}

\begin{proof}
Expand $f_i$ in an orthonormal eigenbasis of the mean-zero subspace. If
$a^2$ is the squared coefficient associated with eigenvalue
$\rho\in[0,1)$, its contribution to the variance is
\[
 \frac{a^2}{T^2}
 \left(T+2\sum_{s=1}^{T-1}(T-s)\rho^s\right)
 \le\frac{2a^2}{T(1-\rho)}.
\]
Summing the contributions gives \eqref{eq:app-coordinate-variance} by the
spectral definition of $H_i$.
\end{proof}

\section{An Explicit Comparison Green Function}
\label{app:profile}

Relabel the positive probabilities so that
$p_1\ge\cdots\ge p_m>0$. Define the graph Laplacian $L_0$ on $[m]$ whose
edge conductances are
\begin{equation}
 w^0_{ij}=\frac{\min\{p_i,p_j\}}{n}.
 \label{eq:app-min-graph}
\end{equation}
For $k\in[m]$, define
\begin{equation}
 c_k=kp_k+\sum_{j>k}p_j
     =\sum_{j=1}^m\min\{p_j,p_k\}.
 \label{eq:app-capped}
\end{equation}

\begin{lemma}[Helmert diagonalization]
\label{lem:app-helmert}
For $k=2,\ldots,m$, let
\begin{equation}
 u_k=\frac{(\underbrace{1,\ldots,1}_{k-1},-(k-1),0,\ldots,0)}
 {\sqrt{k(k-1)}}.
 \label{eq:app-helmert}
\end{equation}
The vectors $(u_k)_{k=2}^m$ are an orthonormal basis of the Euclidean
zero-sum subspace and
\begin{equation}
 L_0u_k=\frac{c_k}{n}u_k.
 \label{eq:app-eigenvalue}
\end{equation}
\end{lemma}

\begin{proof}
Orthonormality is the standard Helmert calculation. Fix $k$. Every coordinate
larger than $k$ sees equal conductance $p_j/n$ to all first $k$ coordinates,
whose entries in $u_k$ sum to zero. Its image is therefore zero. A coordinate
$i<k$ differs only from coordinate $k$ among the first $k$, contributing
$kp_k/n$. Every coordinate $j>k$ contributes $p_j/n$. Thus the $i$th image
is $(c_k/n)u_k(i)$. The same sum at coordinate $k$ is multiplied by
$-(k-1)$, which proves \eqref{eq:app-eigenvalue}.
\end{proof}

Let $D=\operatorname{diag}(p)$. The graph Laplacian associated with
\eqref{eq:app-dirichlet} is $L=D(I-P)$. For the centered indicator $f_i$,
\begin{equation}
 Df_i=p_i(e_i-p)=:b_i.
 \label{eq:app-source}
\end{equation}
Therefore
\begin{equation}
 H_i=b_i^\top L^\dagger b_i.
 \label{eq:app-electrical}
\end{equation}
The exact conductance of $L$ satisfies
\begin{equation}
 \frac{p_ip_j}{(n-1)(p_i+p_j)}
 \ge\frac{\min\{p_i,p_j\}}{2(n-1)}
 \ge\frac12 w^0_{ij}.
 \label{eq:app-loewner}
\end{equation}
It follows by the variational characterization of inverse Laplacians that
\begin{equation}
 H_i\le2b_i^\top L_0^\dagger b_i.
 \label{eq:app-loewner-inverse}
\end{equation}

\begin{lemma}[Exact diagonal formula]
\label{lem:app-diagonal}
Define
\begin{align}
 A_i={}&\sum_{k=2}^{i-1}\frac{(1-c_k)^2}{k(k-1)c_k}
 +\one\{i\ge2\}\frac{(i-c_i)^2}{i(i-1)c_i}\notag\\
 &+\sum_{k=i+1}^{m}\frac{c_k}{k(k-1)}.
 \label{eq:app-Ai}
\end{align}
Then
\begin{equation}
 b_i^\top L_0^\dagger b_i=np_i^2A_i,
 \qquad H_i\le2np_i^2A_i.
 \label{eq:app-diagonal}
\end{equation}
\end{lemma}

\begin{proof}
First observe from \eqref{eq:app-capped} and \eqref{eq:app-helmert} that
\begin{equation}
 u_k^\top p=\frac{1-c_k}{\sqrt{k(k-1)}}.
 \label{eq:app-projected-mean}
\end{equation}
Since $b_i=p_i(e_i-p)$, its projection onto $u_k$ is $p_i$ times
\begin{equation}
 u_k(i)-u_k^\top p=
 \begin{cases}
 -(1-c_k)/\sqrt{k(k-1)}, & k<i,\\
 -(i-c_i)/\sqrt{i(i-1)}, & k=i,\\
 c_k/\sqrt{k(k-1)}, & k>i.
 \end{cases}
 \label{eq:app-projections}
\end{equation}
Divide the squared projections by the eigenvalues $c_k/n$ from
Lemma~\ref{lem:app-helmert} and sum over $k$. This gives the equality in
\eqref{eq:app-diagonal}. The inequality follows from
\eqref{eq:app-loewner-inverse}.
\end{proof}

\section{The Profile Sum Telescopes}
\label{app:telescope}

Write $A_i=L_i+M_i+R_i$ for the three terms in
\eqref{eq:app-Ai}. We prove the aggregate inequality used in the main paper.

\begin{lemma}[Uniform profile bound]
\label{lem:app-profile-sum}
For every nonincreasing positive probability vector,
\begin{equation}
 \sum_{i=1}^mp_i\sqrt{A_i}\le2+2\sqrt2.
 \label{eq:app-profile-sum}
\end{equation}
\end{lemma}

\begin{proof}
By Jensen's inequality and a change in summation order,
\begin{align}
 \sum_i p_i\sqrt{L_i}
 &\le\sqrt{\sum_i p_iL_i}\\
 &=\sqrt{\sum_{k=2}^{m-1}
 \frac{(1-c_k)^2}{k(k-1)c_k}\sum_{i>k}p_i}\\
 &\le\sqrt{\sum_{k=2}^{m-1}\frac1{k(k-1)}}\le1.
 \label{eq:app-left}
\end{align}
The last line uses $\sum_{i>k}p_i\le c_k$ and $1-c_k\le1$. Similarly,
\begin{align}
 \sum_i p_i\sqrt{R_i}
 &\le\sqrt{\sum_i p_iR_i}\\
 &=\sqrt{\sum_{k=2}^{m}\frac{c_k}{k(k-1)}
 \sum_{i<k}p_i}\\
 &\le\sqrt{\sum_{k=2}^{m}\frac1{k(k-1)}}\le1.
 \label{eq:app-right}
\end{align}

For the middle term, define $S_i=\sum_{j\ge i}p_j$. Since
$c_i=ip_i+\sum_{j>i}p_j\ge S_i$ and $i/(i-1)\le2$,
\begin{align}
 \sum_i p_i\sqrt{M_i}
 &\le\sqrt2\sum_i\frac{p_i}{\sqrt{c_i}}
 \le\sqrt2\sum_i\frac{S_i-S_{i+1}}{\sqrt{S_i}}.
 \label{eq:app-middle1}
\end{align}
For $0\le b\le a$,
\begin{equation}
 \frac{a-b}{\sqrt a}\le2(\sqrt a-\sqrt b).
 \label{eq:app-sqrt-telescope}
\end{equation}
Apply this with $a=S_i,b=S_{i+1}$ and telescope to obtain
\begin{equation}
 \sum_i p_i\sqrt{M_i}\le2\sqrt2.
 \label{eq:app-middle2}
\end{equation}
Finally, $\sqrt{L_i+M_i+R_i}\le\sqrt{L_i}+\sqrt{M_i}+\sqrt{R_i}$.
Combining \eqref{eq:app-left}, \eqref{eq:app-right}, and
\eqref{eq:app-middle2} proves the claim.
\end{proof}

\section{Stationary Occupation Risk}
\label{app:stationary}

\begin{theorem}[Stationary risk]
\label{thm:app-risk}
For a stationary length-$T$ winner-chain trajectory with empirical law
$\widehat p_T$,
\begin{equation}
 \E\TV(\widehat p_T,p)
 \le(2+2\sqrt2)\sqrt{\frac nT}.
 \label{eq:app-risk}
\end{equation}
\end{theorem}

\begin{proof}
If the support has size one, the claim is immediate. Otherwise,
Lemmas~\ref{lem:app-coordinate-variance} and \ref{lem:app-diagonal} give
\[
 \operatorname{Var}(\widehat p_T(i))
 \le\frac{4np_i^2A_i}{T}.
\]
Jensen's inequality and Lemma~\ref{lem:app-profile-sum} yield
\begin{align*}
 \E\TV(\widehat p_T,p)
 &\le\frac12\sum_i\E|\widehat p_T(i)-p_i|\\
 &\le\frac12\sum_i\sqrt{\operatorname{Var}(\widehat p_T(i))}\\
 &\le\sqrt{\frac nT}\sum_i p_i\sqrt{A_i}
 \le(2+2\sqrt2)\sqrt{\frac nT}.
\end{align*}
Zero-mass coordinates are never visited under stationarity and add zero.
\end{proof}

\section{Arbitrary Starts and High Confidence}
\label{app:amplification}

Set $C_0=2+2\sqrt2$ and
\begin{equation}
 B=\left\lceil(n-1)\log\frac{48}{\varepsilon}\right\rceil,
 \qquad
 T=\left\lceil48^2C_0^2\frac{n}{\varepsilon^2}\right\rceil.
 \label{eq:app-BT}
\end{equation}
Start the winner chain from an arbitrary state and form its empirical law
$\widehat p$ over the $T$ states after burn-in.

\begin{lemma}[One base estimate]
\label{lem:app-base}
The estimate satisfies
\begin{equation}
 \Prb\{\TV(\widehat p,p)>\varepsilon/4\}\le1/6.
 \label{eq:app-base}
\end{equation}
\end{lemma}

\begin{proof}
By \eqref{eq:app-mixing}, the state after burn-in is within
$\varepsilon/48$ TV of stationarity. Generate the entire subsequent path by
applying one common path kernel to either initial law. Data processing shows
that the arbitrary-start and stationary path laws remain within
$\varepsilon/48$ TV. Since TV loss lies in $[0,1]$, their expected losses
differ by at most $\varepsilon/48$. Theorem~\ref{thm:app-risk} and the choice
of $T$ give
\[
 \E\TV(\widehat p,p)\le\varepsilon/48+\varepsilon/48
 =\varepsilon/24.
\]
Markov's inequality proves \eqref{eq:app-base}.
\end{proof}

Repeat this construction independently $R$ times and let $\widehat p_r$ be
the candidates. Define an $L_1$ geometric median
\begin{equation}
 \widetilde p\in\arg\min_{q\in\Delta_n}
 \sum_{r=1}^R\TV(q,\widehat p_r).
 \label{eq:app-median}
\end{equation}
Existence follows from compactness. The optimization is a linear program
after introducing one absolute-value slack variable per candidate and
coordinate.

\begin{lemma}[Metric median]
\label{lem:app-median}
In any metric space, suppose $G>R/2$ of the points $z_1,\ldots,z_R$ lie
within distance $r$ of $p$. Every minimizer $z$ of the sum of distances to
the points satisfies
\begin{equation}
 d(z,p)\le\frac{2G}{2G-R}r.
 \label{eq:app-median-bound}
\end{equation}
\end{lemma}

\begin{proof}
Let $\mathcal G$ denote the good indices and $d=d(z,p)$. For a good point,
\[
 d(z,z_r)-d(p,z_r)\ge d-2r.
\]
For a bad point, the triangle inequality gives
$d(z,z_r)-d(p,z_r)\ge-d$. Since the sum at $z$ is no larger than the sum at
$p$,
\[
 0\ge G(d-2r)-(R-G)d=(2G-R)d-2Gr.
\]
Rearranging proves the result.
\end{proof}

\begin{theorem}[High-confidence upper bound]
\label{thm:app-upper}
If $R=\lceil18\log(1/\delta)\rceil$, then
\begin{equation}
 \Prb\{\TV(\widetilde p,p)>\varepsilon\}\le\delta.
 \label{eq:app-upper}
\end{equation}
The total number of pair queries is
$O(n\log(1/\delta)/\varepsilon^2)$.
\end{theorem}

\begin{proof}
Every candidate is good with probability at least $5/6$ by
Lemma~\ref{lem:app-base}. Hoeffding's inequality gives
\[
 \Prb\{G<2R/3\}\le e^{-R/18}\le\delta.
\]
On the complementary event, apply Lemma~\ref{lem:app-median} with
$r=\varepsilon/4$. The multiplicative factor in
\eqref{eq:app-median-bound} is at most four, so the output error is at most
$\varepsilon$.

Each repetition uses $B+T$ queries. For $0<\varepsilon\le1/16$,
$\log(48/\varepsilon)=O(1/\varepsilon^2)$, hence
$B+T=O(n/\varepsilon^2)$. Multiplying by $R$ proves the query bound.
\end{proof}

\section{Adaptive Confidence Lower Bound}
\label{app:lower}

Let $n\ge4$, $0<\varepsilon\le1/16$, and define
\begin{equation}
 p_a(1)=a,\qquad p_a(j)=\frac{1-a}{n-1},\quad j\ge2.
 \label{eq:app-hard-family}
\end{equation}
Take $a_-=1/2-2\varepsilon$ and $a_+=1/2+2\varepsilon$. Both lie in
$[1/4,3/4]$, and
\begin{equation}
 \TV(p_{a_-},p_{a_+})=|a_--a_+|=4\varepsilon.
 \label{eq:app-separation}
\end{equation}

If a queried pair excludes label $1$, its response is uniform under both
hypotheses. If the pair is $\{1,j\}$, write one for a light-label response.
The Bernoulli parameter is
\begin{equation}
 r(a)=\frac{1-a}{1+(n-2)a}.
 \label{eq:app-r}
\end{equation}
For $a\in[1/4,3/4]$ and $n\ge4$,
\begin{equation}
 \frac1{3n}\le r(a)\le\frac3n,
 \qquad
 |r'(a)|=\frac{n-1}{(1+(n-2)a)^2}\le\frac{16}{n}.
 \label{eq:app-r-bounds}
\end{equation}
Also $1-r(a)\ge1/4$. By the mean value theorem and the Bernoulli inequality
\begin{equation}
 D_{\rm KL}(\operatorname{Bern}(x)\|
 \operatorname{Bern}(y))
 \le\frac{(x-y)^2}{y(1-y)},
 \label{eq:app-bern-kl}
\end{equation}
every possible pair query has conditional KL divergence at most
\begin{equation}
 C_1\frac{\varepsilon^2}{n}
 \label{eq:app-per-query}
\end{equation}
between the two response laws, for a universal $C_1$.

\begin{theorem}[Adaptive lower bound]
\label{thm:app-lower}
Any possibly randomized and adaptive learner using a fixed budget $Q$ and
satisfying
\begin{equation}
 \sup_{p\in\Delta_n}
 \Prb_p\{\TV(\widehat p,p)>\varepsilon\}\le\delta
 \label{eq:app-learning-guarantee}
\end{equation}
must have
\begin{equation}
 Q\ge c\frac{n\log(1/\delta)}{\varepsilon^2}
 \label{eq:app-lower}
\end{equation}
for $0<\delta\le1/8$ and a universal $c>0$.
\end{theorem}

\begin{proof}
Include the learner's private random seed in its transcript. Conditional on
every past transcript, its next pair is fixed. The conditional response laws
then obey \eqref{eq:app-per-query}, regardless of which pair was selected.
The KL chain rule gives
\begin{equation}
 D_{\rm KL}(\mathsf P_-^Q\|\mathsf P_+^Q)
 \le C_1Q\frac{\varepsilon^2}{n},
 \label{eq:app-transcript-kl}
\end{equation}
where $\mathsf P_-^Q,\mathsf P_+^Q$ are the complete transcript laws.

Given the output $\widehat p$, test the two hypotheses by choosing the closer
of $p_{a_-}$ and $p_{a_+}$ in TV. By
\eqref{eq:app-separation}, every estimate within $\varepsilon$ of the truth
is strictly closer to that truth. Thus the two testing errors are at most
$\delta$. The Bretagnolle--Huber inequality
\citep{tsybakov2009introduction} gives
\begin{equation}
 2\delta\ge\frac12
 \exp\{-D_{\rm KL}(\mathsf P_-^Q\|\mathsf P_+^Q)\}.
 \label{eq:app-BH}
\end{equation}
Combining \eqref{eq:app-transcript-kl} and \eqref{eq:app-BH} yields
\[
 Q\ge\frac{n}{C_1\varepsilon^2}\log\frac1{4\delta}.
\]
For $\delta\le1/8$, this is \eqref{eq:app-lower} after changing the
universal constant.
\end{proof}

Theorem~\ref{thm:app-upper} and Theorem~\ref{thm:app-lower} together prove
the minimax theorem in the main paper.

\section{A Bounded-Menu Confidence Lower Bound}
\label{app:bounded-menu}

The same family quantifies how menu size can reduce the confidence
obstruction. Suppose every allowed conditional query has cardinality at most
$k$, where $2\le k\le n$.

\begin{proposition}[Bounded-menu lower bound]
\label{prop:app-bounded-menu}
Under the parameter ranges of Theorem~\ref{thm:app-lower}, every learner from
menus of size at most $k$ requires
\begin{equation}
 Q=\Omega\!\left(
 \frac{n}{k}\frac{\log(1/\delta)}{\varepsilon^2}
 \right).
 \label{eq:app-bounded-menu}
\end{equation}
\end{proposition}

\begin{proof}
Use the two distributions in \eqref{eq:app-hard-family}. A menu excluding
label $1$ is uniform under both hypotheses and carries no information. If a
menu contains label $1$ and $s\le k-1$ light labels, the identity of a light
winner is uniform conditional on a light outcome. The only informative
statistic is therefore Bernoulli with light-outcome probability
\begin{equation}
 r_s(a)=\frac{s(1-a)}{(n-1)a+s(1-a)}.
 \label{eq:app-rs}
\end{equation}
Uniformly for $a\in[1/4,3/4]$,
$r_s(a)=\Theta(s/n)$ and
\begin{equation}
 |r_s'(a)|
 =\frac{s(n-1)}{((n-1)a+s(1-a))^2}
 =O(s/n).
 \label{eq:app-rs-derivative}
\end{equation}
Equation \eqref{eq:app-bern-kl} now bounds the one-query KL by
$Ck\varepsilon^2/n$. The adaptive transcript and testing argument from
Theorem~\ref{thm:app-lower} gives \eqref{eq:app-bounded-menu}.
\end{proof}

Ordinary distribution-learning hardness also gives
$\Omega(n/\varepsilon^2)$ even when the full domain is queryable
\citep{canonne2020short}. Hence every bounded-menu learner has lower bound
\[
 \Omega\!\left(\frac{n+(n/k)\log(1/\delta)}{\varepsilon^2}\right).
\]
We do not claim a matching upper bound when $k>2$.

\section{Mechanical Verification}
\label{app:verification}

The proofs are analytic. The script
\texttt{verify\_trajectory\_theorem.py} performs independent floating-point
checks designed to catch indexing, scaling, and inequality-direction errors.
It does not supply evidence needed by any theorem.

Across 1,014 seeded mass profiles in dimensions $2$ through $40$, including
uniform, one-heavy, geometric, and Dirichlet profiles, the script checked:
\begin{enumerate}
\item the winner kernel's symmetrized eigenvalues lie in $[0,1]$,
\item the Helmert formula \eqref{eq:app-diagonal} agrees with a direct
pseudoinverse of $L_0$,
\item the exact Green function is bounded by twice the comparison Green
function coordinatewise,
\item each of the left, middle, and right profile sums obeys its analytic
bound,
\item the exact finite-$T$ stationary covariance for
$T\in\{1,2,5,20,100\}$ obeys
Lemma~\ref{lem:app-coordinate-variance}, and
\item the normalized one-query KL in the hard family stays below the loose
analytic constant over $4\le n\le256$ and 80 logarithmically spaced
accuracies.
\end{enumerate}
The largest observed normalized KL,
$nD_{\rm KL}/\varepsilon^2$, was $161.847400$. The terminal output was
\texttt{all trajectory-theorem checks passed}.

An additional stress script sampled highly imbalanced simplex points through
dimension $1024$. The observed asymptotic $L_1$ risk divided by
$\sqrt{n/T}$ remained below $2.28$, while the theorem uses the conservative
TV constant $2+2\sqrt2$. These searches motivated no step of the proof.

\section{Source and Novelty Audit}
\label{app:sources}

We read the full 34-page paper of \citet{kleinberg2026providers}. Its generic
complete-family learner runs a global winner chain, burns in, keeps one
terminal state, and independently restarts. Its hierarchical learner keeps
local trajectories but estimates one binary tree split at a time. Its
pairwise corollary is
$O(n\log^2n\log(n/\delta)/\varepsilon^2)$, and its conclusion explicitly
leaves the polylogarithmic gaps open. It contains neither the empirical-law
risk theorem nor the multiplicative-confidence lower bound.

We also read the full papers on accelerated spectral ranking
\citep{agarwal2018accelerated}, graph-resistance estimation
\citep{hendrickx2019resistance}, and the asymptotic BTL minimax rate
\citep{hendrickx2020minimax}, together with the full paper and supplement of
\citet{shah2015estimation}. They study fixed comparison designs and parameter
losses. Their useful uniform guarantees assume positive weights and bounded
parameter ranges or enter a many-comparisons-per-edge regime. None analyzes
the query-driven winner trajectory used here.

The full exact-sampling paper of \citet{fotakis2022perfect} uses coupling
from the past for an exogenous pair design. Its distribution-free procedure
first estimates all weights to relative accuracy under a learnability
condition, then modifies the chain and performs exact correction. Its goal
and complexity are distinct from minimax labeled TV learning.

For statistical context, we read the complete ordinary distribution-learning
note of \citet{canonne2020short}, the full Markov Bernstein paper of
\citet{jiang2018bernstein}, and the classical biased-sampling sources
\citep{vardi1985empirical,gill1988large}. The first gives the additive
ordinary-sampling confidence benchmark. Generic spectral-gap concentration
from the second incurs the relaxation-time factor avoided by the capped-mass
profile. The biased-sampling results are asymptotic and fixed-design.

A live arXiv API, arXiv web, and PMLR search through July 30, 2026 combined
pairwise conditional sampling, PCOND, Bradley--Terry--Luce distribution
learning, stationary-law estimation, empirical trajectories, total
variation, and confidence. It returned the sources above but no theorem with
the rate in the main paper, no capped-mass Helmert analysis, and no matching
high-confidence lower bound. The closest source,
\citet{kleinberg2026providers}, was posted on July 27, 2026 and is cited as
the direct starting point rather than presented as concurrent work.

\bibliography{references}

\end{document}
